% Reader copy: keep the opening word with its display; frozen source unchanged.
% Live-book statement adaptation of the arbitrary-ground-field nonreduced
% addendum.  Date: 28 August 2026.

\begin{proposition}
\label{prop:p21p3-curve-rigidity-nonreduced}
Let \(K\) be a complete
nontrivially valued non-Archimedean field and let \(X\) be a good separated
strict \(K\)-analytic space, pure of dimension one and without boundary.
The space \(X\) is not assumed reduced.  Write

\[
             \iota:X_{\mathrm{red}}\hookrightarrow X
\]

for the canonical nilpotent closed immersion.  Let \((X_i)_i\) be the
irreducible components of \(X_{\mathrm{red}}\), assume that their
normalizations are rig-smooth, and write

\[
 \nu_i:\widetilde X_i\longrightarrow X_i
       \hookrightarrow X_{\mathrm{red}}\xrightarrow{\iota}X,
 \qquad
 \delta_i:=(\nu_i)_*\delta_{\widetilde X_i}.
\]

Let \(m_i\in\mathbf Z_{>0}\) be the coefficient of \([X_i]\) in the
fundamental cycle \([X]\).  Then the following assertions hold.

\begin{enumerate}
\item\label{itm:p21p3-nr:calculus}
For every \((p,q)\), pullback by \(\iota\) identifies the sheaves of
ordinary CLD-smooth forms on \(X\) and \(X_{\mathrm{red}}\), compatibly
with wedge product, \(J\), \(\mathrm d'\), \(\mathrm d''\), support, and all
three positivity cones.  On compactly supported test forms this is an
isomorphism of topological vector spaces.  Its transpose identifies the
CLD current sheaves, their weak topologies, positivity, and
\(\mathrm d'\)- and \(\mathrm d''\)-closedness.  It identifies multiplication
by smooth functions for all currents.  If \(S\) is positive and \(f\) is a
locally bounded Borel function, it also identifies the CLD product \(fS\);
by linearity the same holds for currents locally expressible as differences
of positive currents.  No multiplication by a Borel function is asserted
for an arbitrary order-zero current.  Under the common underlying
topological space, the CLD-psh and continuous CLD-psh sheaves on \(X\) and
\(X_{\mathrm{red}}\) coincide.  The analogous pullback identification
holds for \(G\)-forms.

\item\label{itm:p21p3-nr:integration}
The integration current is not, in general, invariant under reduction.
Precisely,

\begin{equation}
\label{eq:p21p3-nr-fundamental-cycle}
        \delta_X=\sum_i m_i\delta_i,
        \qquad
        \iota_*\delta_{X_{\mathrm{red}}}=\sum_i\delta_i.
\end{equation}

Both sums are locally finite.  In particular,
\(\delta_X=\iota_*\delta_{X_{\mathrm{red}}}\) if and only if \(X\) is
generically reduced along every \(X_i\).

\item\label{itm:p21p3-nr:classification}
Every positive \(\mathrm d'\)- and \(\mathrm d''\)-closed CLD current \(T\)
of bidegree \((0,0)\) on \(X\) has a unique expression

\begin{equation}
\label{eq:p21p3-nr-classification}
                     T=\sum_i c_i\delta_i,
                     \qquad c_i\in\mathbf R_{\geq0}.
\end{equation}

The sum is locally finite.  If \(X_{\mathrm{red}}\) is connected and
rig-smooth, there is one generic multiplicity \(m\), and every such current
has the form \(T=c\delta_X\) for a unique \(c\geq0\).

\item\label{itm:p21p3-nr:bedford-taylor}
Let \(W\subseteq X\) be open and let \(U,V\) be continuous CLD-psh
functions on \(W\).  Identify \(W\) with its inverse image in
\(X_{\mathrm{red}}\), and, for every continuous function \(F\) on \(W\),
put

\[
 [F]_X:=F(\delta_X|_W),\qquad
 [F]_{X_{\mathrm{red}}}:=F(\delta_{X_{\mathrm{red}}}|_W).
\]

With the identifications in \cref{itm:p21p3-nr:calculus}, one has the
exact equivalence

\begin{equation}
\label{eq:p21p3-nr-curvature-equivalence}
 \mathrm d'\mathrm d''[U]_X=\mathrm d'\mathrm d''[V]_X
 \quad\Longleftrightarrow\quad
 \mathrm d'\mathrm d''[U]_{X_{\mathrm{red}}}
   =\mathrm d'\mathrm d''[V]_{X_{\mathrm{red}}}.
\end{equation}

If these equivalent conditions hold, then every \(T\) in
\cref{itm:p21p3-nr:classification} satisfies

\begin{equation}
\label{eq:p21p3-nr-bt-conclusion}
       \mathrm d'\mathrm d''\bigl((U-V)T|_W\bigr)=0.
\end{equation}

Consequently every positive closed CLD current of bidegree \((0,0)\) on
\(X\) has the Bedford--Taylor independence property for continuous psh
pairs.  Under the current identification in
\cref{itm:p21p3-nr:calculus}, this continuous-pair property is invariant
under the nilpotent thickening
\(X_{\mathrm{red}}\hookrightarrow X\).
\end{enumerate}
\end{proposition}

\begin{remark}
\label{rem:p21p3-nr-scope}
No algebraic-closedness hypothesis is made.  Strictness is used
to ensure that the zero-dimensional conductor locus and its inverse image
in every normalization consist of locally finite rigid points; the rigid
zero-atom lemma then replaces the algebraically closed/type-\(1\) input.
The addendum does not treat a trivially valued field, a non-strict or
non-good curve, boundary points, a component with non-rig-smooth
normalization, or analytic dimension at least two.  The last item concerns
only continuous psh pairs.  Since CLD Definition~8.6.2 literally
quantifies over psh functions that may be discontinuous, no assertion of
that stronger predicate is made here.  Nilpotents enter only through the
generic multiplicities \(m_i\); no nilpotence exponent or transverse
vanishing-order estimate is used.
\end{remark}

\begin{samepage}
\begin{remark}
\label{rem:p21p3-nr-double-line}
For
\[
 X=\bigl(\operatorname{Spec}K[T,\epsilon]/(\epsilon^2)\bigr)^{\mathrm{an}},
 \qquad X_{\mathrm{red}}=\mathbf A_K^{1,\mathrm{an}},
\]

the unique generic multiplicity is \(2\).  Hence all smooth forms and all
positive closed currents are identified with those on the analytic line,
whereas
\(\delta_X=2\iota_*\delta_{\mathbf A_K^{1,\mathrm{an}}}\).
Thus generic multiplicity, and not a new nilpotent differential form or
current, is the precise obstruction to reduction-invariance of the
integration current.  The proof uses only the normalizations of the
components of \(X_{\mathrm{red}}\); no normalization of the nonreduced
space is needed.
\end{remark}
\end{samepage}
